\chapter{1977年江苏南通市初试卷}

\section{解答题}

\begin{problem}
化简，计算：

\begin{enumerate}
\item \(\dfrac{-2mn}{m-3n}\sqrt{\dfrac{\left(m+3n\right)^2}{2mn}}-6\)（\(0<m<3n\)）；
\item \(\left(\dfrac14\right)^{-1.5}\cdot\left[\left(\dfrac7{11}\right)^0\right]^{-3}-\left(\dfrac35\right)^{-2}\div\left(2\dfrac79\right)^{\frac12}\).
\end{enumerate}
\end{problem}

\begin{solution}
\begin{enumerate}
\item 由 \(0<m<3n\) 可得 \(n>\dfrac{m}{3}>0\)，所以 \(2mn>0\)，且 \(m+3n>0\). 因而 \(\sqrt{\dfrac{\left(m+3n\right)^2}{2mn}}=\dfrac{m+3n}{\sqrt{2mn}}\). 又因为 \(m-3n<0\)，所以原式为
\[
\dfrac{-2mn}{m-3n}\sqrt{\dfrac{\left(m+3n\right)^2}{2mn}}-6
=\dfrac{2mn}{3n-m}\cdot\dfrac{m+3n}{\sqrt{2mn}}-6
=\dfrac{\left(m+3n\right)\sqrt{2mn}}{3n-m}-6
\]

\item 先将小数指数和带分数化为分数，得 \(1.5=\dfrac32\)，\(2\dfrac79=\dfrac{25}{9}\). 由于 \(\left(\dfrac7{11}\right)^0=1\)，且 \(\sqrt{\dfrac{25}{9}}=\dfrac53\)，所以
\[
\left(\dfrac14\right)^{-1.5}\left[\left(\dfrac7{11}\right)^0\right]^{-3}-\left(\dfrac35\right)^{-2}\div\left(2\dfrac79\right)^{\frac12}
=4^{\frac32}\cdot1^{-3}-\left(\dfrac53\right)^2\div\sqrt{\dfrac{25}{9}}
=8-\dfrac{25}{9}\div\dfrac53
=8-\dfrac{25}{9}\cdot\dfrac35=\dfrac{19}{3}
\]
\end{enumerate}
\end{solution}

\begin{problem}
在实数范围内分解因式：

\begin{enumerate}
\item \(a^4-10a^2+16\)；
\item \(x^2-4xy+4y^2-3x+6y+2\).
\end{enumerate}
\end{problem}

\begin{solution}
\begin{enumerate}
\item 令 \(u=a^2\)，则
\[
a^4-10a^2+16=u^2-10u+16=\left(u-2\right)\left(u-8\right)
\]
代回 \(u=a^2\)，并在实数范围内继续分解，得
\[
a^4-10a^2+16=\left(a^2-2\right)\left(a^2-8\right)
=\left(a-\sqrt2\right)\left(a+\sqrt2\right)\left(a-2\sqrt2\right)\left(a+2\sqrt2\right)
\]

\item 将含有 \(x\)，\(y\) 的项按 \(x-2y\) 重新组合，得
\[
x^2-4xy+4y^2-3x+6y+2
=\left(x-2y\right)^2-3\left(x-2y\right)+2
\]
令 \(v=x-2y\)，则
\[
v^2-3v+2=\left(v-1\right)\left(v-2\right)
\]
代回 \(v=x-2y\)，得
\[
x^2-4xy+4y^2-3x+6y+2=\left(x-2y-1\right)\left(x-2y-2\right)
\]
\end{enumerate}
\end{solution}

\begin{problem}
求 \(x\)：

\begin{enumerate}
\item \(\dfrac{2x+1}{x-3}-\dfrac{5x-11}{x^2-4x+3}=\dfrac3{x-1}\).
\item \(\sqrt{2x-4}-\sqrt{x+5}=1\).
\item \(\log_{x-1}\left(2x^2-3x-1\right)=2\).
\end{enumerate}
\end{problem}

\begin{solution}
\begin{enumerate}
\item 因为分母中出现 \(x-3\) 和 \(x^2-4x+3\)，且 \(x^2-4x+3=\left(x-1\right)\left(x-3\right)\)，所以原方程的定义域为 \(x\ne1\)，\(3\). 在这个条件下，两边同乘 \(\left(x-1\right)\left(x-3\right)\)，得到
\[
\left(2x+1\right)\left(x-1\right)-\left(5x-11\right)=3\left(x-3\right)
\]
展开并移项，得
\(2x^2-6x+10=3x-9\)，\(2x^2-9x+19=0\).
这个二次方程的判别式为 \(\Delta=(-9)^2-4\times2\times19=81-152=-71<0\)，所以它没有实数根，原方程无实数解.

\item 两个根式都有意义，首先要求 \(2x-4\geqslant0\)，即 \(x\geqslant2\). 移项得 \(\sqrt{2x-4}=1+\sqrt{x+5}\). 两边非负，可以平方，得 \(2x-4=1+x+5+2\sqrt{x+5}\)，即 \(x-10=2\sqrt{x+5}\). 右端非负，所以还必须有 \(x\geqslant10\). 在此条件下再次平方，得
\(\left(x-10\right)^2=4\left(x+5\right)\)，\(x^2-24x+80=0\).
解得 \(x=4\) 或 \(x=20\). 结合 \(x\geqslant10\) 只保留 \(x=20\)，并检验 \(\sqrt{2\times20-4}-\sqrt{20+5}=6-5=1\)，因此原方程的解为 \(x=20\).

\item 对数底数的条件为 \(x-1>0\) 且 \(x-1\ne1\)，即 \(x>1\) 且 \(x\ne2\). 真数还必须为正. 方程 \(2x^2-3x-1=0\) 的根为 \(x=\dfrac{3\pm\sqrt{17}}{4}\). 在 \(x>1\) 的范围内，真数为正还要求 \(x>\dfrac{3+\sqrt{17}}4\). 在这些定义域条件下，由对数定义，原方程等价于 \(2x^2-3x-1=\left(x-1\right)^2\). 整理得 \(x^2-x-2=0\)，即 \(\left(x-2\right)\left(x+1\right)=0\)，所以候选值为 \(x=2\)，\(-1\). 但 \(x=2\) 时底数为 \(1\)，\(x=-1\) 时底数为负数，二者均不在定义域内，因此原方程无实数解.
\end{enumerate}
\end{solution}

\begin{problem}
求下列函数的定义域：

\begin{enumerate}
\item \(y=\dfrac1{1-\sqrt{x-1}}\)；
\item \(y=\lg\left(\tan x\right)\).
\end{enumerate}
\end{problem}

\begin{solution}
\begin{enumerate}
\item 根式有意义要求 \(x-1\geqslant0\)，即 \(x\geqslant1\). 同时分母不能为零，故 \(1-\sqrt{x-1}\ne0\)，这等价于 \(\sqrt{x-1}\ne1\)，即 \(x\ne2\). 因此定义域为 \(x\geqslant1\) 且 \(x\ne2\)，也可写为 \([1,2)\cup(2,+\infty)\).

\item \(\lg\) 的真数必须为正，因此要求 \(\tan x>0\). 这个条件同时排除了正切函数无定义的点. 正切函数在每个周期内为正的区间是
\(\left(k\pi,k\pi+\dfrac\pi2\right)\)，\(k\in\Z\).
所以定义域为
\[
\bigcup_{k\in\Z}\left(k\pi,k\pi+\dfrac\pi2\right)
\]
\end{enumerate}
\end{solution}

\begin{problem}
不查表：

\begin{enumerate}
\item 求 \(4^{1-\log_2\dfrac25+\log_2\sqrt2}\).
\item 解三角形：已知 \(\triangle ABC\) 中，\(\angle B=120^\circ\)，\(\angle C=15^\circ\)，\(b=100\)，求 \(c\)（可用根式表示结果）.
\item 将直角坐标方程 \(x^2-y^2=a^2\) 化成极坐标方程.
\end{enumerate}
\end{problem}

\begin{solution}
\begin{enumerate}
\item 由对数的商和幂的运算性质，\(\log_2\dfrac25=\log_2 2-\log_2 5=1-\log_2 5\)，且 \(\log_2\sqrt2=\dfrac12\). 因此原式的指数为
\[
1-\log_2\dfrac25+\log_2\sqrt2
=1-\left(1-\log_2 5\right)+\dfrac12
=\dfrac12+\log_2 5
\]
于是
\[
4^{1-\log_2\frac25+\log_2\sqrt2}=\left(2^2\right)^{\frac12+\log_2 5}=2^{1+2\log_2 5}=2\cdot5^2=50
\]

\item 由三角形内角和，\(\angle A=180^\circ-120^\circ-15^\circ=45^\circ\). 由正弦定理，\(\dfrac{c}{\sin15^\circ}=\dfrac{b}{\sin120^\circ}\). 又因为
\[
\sin15^\circ=\sin\left(45^\circ-30^\circ\right)=\sin45^\circ\cos30^\circ-\cos45^\circ\sin30^\circ=\dfrac{\sqrt6-\sqrt2}{4}
\]
同时 \(\sin120^\circ=\sin60^\circ=\dfrac{\sqrt3}{2}\)，所以
\[
c=100\cdot\dfrac{\sin15^\circ}{\sin120^\circ}=100\cdot\dfrac{\frac{\sqrt6-\sqrt2}{4}}{\frac{\sqrt3}{2}}=\dfrac{50\left(\sqrt6-\sqrt2\right)}{\sqrt3}=\dfrac{50}{3}\left(3\sqrt2-\sqrt6\right)
\]

\item 设动点的极坐标为 \(\left(r,\theta\right)\)，则 \(x=r\cos\theta\)，\(y=r\sin\theta\). 代入 \(x^2-y^2=a^2\)，得 \(r^2\cos^2\theta-r^2\sin^2\theta=a^2\). 利用 \(\cos^2\theta-\sin^2\theta=\cos2\theta\)，得到极坐标方程
\[
r^2\cos2\theta=a^2
\]
\end{enumerate}
\end{solution}

\begin{problem}
\(\triangle ABC\) 的 \(\angle A\) 的平分线交 \(BC\) 于 \(D\)，延长 \(AD\) 交三角形的外接圆于 \(E\)，求证：\(AB\cdot AC=AD\cdot AE\).
\end{problem}

\begin{solution}
设 \(AB=m\)，\(AC=n\)，\(BD=p\)，\(DC=q\). 由于 \(AD\) 是角平分线，角平分线定理给出 \(\dfrac{p}{q}=\dfrac{m}{n}\)，即 \(np=mq\).
在三角形 \(ABC\) 中对截线 \(AD\) 使用斯特瓦特定理，得
\[
n^2p+m^2q=\left(p+q\right)\left(AD^2+pq\right)
\]
由 \(np=mq\)，斯特瓦特定理左端为
\[
n^2p+m^2q=n(np)+m(mq)=mnq+mnp=mn\left(p+q\right)
\]
因此 \(mn\left(p+q\right)=\left(p+q\right)\left(AD^2+pq\right)\). 因为 \(p+q=BC>0\)，约去 \(p+q\)，得到 \(AD^2=mn-pq=AB\cdot AC-BD\cdot DC\). 又因为 \(A\)，\(B\)，\(C\)，\(E\) 共圆，且弦 \(AE\) 与弦 \(BC\) 交于 \(D\)，由相交弦定理 \(AD\cdot DE=BD\cdot DC\). 由题意，点的顺序为 \(A\)，\(D\)，\(E\)，所以 \(AE=AD+DE\). 于是
\[
AB\cdot AC=AD^2+BD\cdot DC=AD^2+AD\cdot DE=AD\left(AD+DE\right)=AD\cdot AE
\]
故 \(AB\cdot AC=AD\cdot AE\).
\end{solution}

\begin{problem}
已知正三棱锥 \(P-ABC\) 的底面边长是 \(6\mathrm{~cm}\)，侧面与底面所成的角是 \(60^\circ\)，求它的体积（结果可用根式表示）.
\end{problem}

\begin{solution}
设正三角形 \(ABC\) 的中心为 \(O\)，\(PO\) 为棱锥的高，\(M\) 为 \(BC\) 的中点. 正三棱锥的高垂直于底面，且侧面三角形 \(PBC\) 为等腰三角形，因此 \(PO\perp\text{平面 }ABC\)，\(PM\perp BC\)，\(OM\perp BC\).
所以侧面 \(PBC\) 与底面 \(ABC\) 所成的二面角等于 \(\angle PMO=60^\circ\)，且三角形 \(POM\) 在 \(O\) 处为直角三角形.

正三角形边长为 \(6\mathrm{~cm}\)，其中心到边的距离为内切圆半径，故 \(OM=\dfrac{6\sqrt3}{6}=\sqrt3\mathrm{~cm}\). 底面积为 \(S_{\triangle ABC}=\dfrac{6^2\sqrt3}{4}=9\sqrt3\mathrm{~cm}^2\). 在直角三角形 \(POM\) 中，\(\tan60^\circ=\dfrac{PO}{OM}\)，所以 \(PO=OM\tan60^\circ=\sqrt3\cdot\sqrt3=3\mathrm{~cm}\).
因此棱锥体积为
\[
V=\dfrac13S_{\triangle ABC}\cdot PO=\dfrac13\cdot9\sqrt3\cdot3=9\sqrt3\mathrm{~cm}^3
\]
\end{solution}

\begin{problem}
已知 \(a\)，\(b\)，\(c\) 为三角形的三条边，求证：\(b^2x^2+\left(b^2+c^2-a^2\right)x+c^2=0\) 无实数根.
\end{problem}

\begin{solution}
因为 \(a\)，\(b\)，\(c\) 是三角形的三条边，所以满足严格三角不等式 \(\left|b-c\right|<a<b+c\). 三边均为正数，平方后得到
\[
\left(b-c\right)^2<a^2<\left(b+c\right)^2
\]
将两端展开并移项，得 \(-2bc<b^2+c^2-a^2<2bc\). 由于 \(b\)，\(c>0\)，上式说明 \(\left|b^2+c^2-a^2\right|<2bc\)，从而
\[
\left(b^2+c^2-a^2\right)^2<4b^2c^2
\]
原方程的二次项系数为 \(b^2>0\)，其判别式为 \(\Delta=\left(b^2+c^2-a^2\right)^2-4b^2c^2<0\).
因此原方程无实数根.
\end{solution}

\begin{problem}
方程 \(x^2-\left(\tan\theta+\cot\theta\right)x+1=0\) 有一根为 \(2+\sqrt3\)，求 \(\sin 2\theta\).
\end{problem}

\begin{solution}
设方程的另一根为 \(x_2\). 由韦达定理，两根之积为 \(1\)，所以 \(\left(2+\sqrt3\right)x_2=1\)，即 \(x_2=\dfrac{1}{2+\sqrt3}=2-\sqrt3\). 再由两根之和，得到 \(\tan\theta+\cot\theta=\left(2+\sqrt3\right)+\left(2-\sqrt3\right)=4\). 因为题目中的 \(\tan\theta\) 和 \(\cot\theta\) 均有意义，所以 \(\sin\theta\cos\theta\ne0\). 因而
\[
\tan\theta+\cot\theta=\dfrac{\sin^2\theta+\cos^2\theta}{\sin\theta\cos\theta}=\dfrac{1}{\sin\theta\cos\theta}=\dfrac{2}{\sin2\theta}
\]
所以 \(\dfrac{2}{\sin2\theta}=4\)，从而 \(\sin2\theta=\dfrac12\).
\end{solution}

\begin{problem}
汽车与自行车分别从相距 \(75\mathrm{~km}\) 的 \(A\)，\(B\) 两地同时出发（\(A\) 地在 \(B\) 地的正东方向），汽车从 \(A\) 地向正西方向行驶，自行车从 \(B\) 地向正南方向行驶，已知汽车和自行车的速度分别为 \(30\mathrm{~km/h}\) 和 \(15\mathrm{~km/h}\)，问出发后多少小时两车相距最短？
\end{problem}

\begin{solution}
取 \(B\) 为坐标原点，正东方向为 \(x\) 轴正向，正北方向为 \(y\) 轴正向，则 \(A=\left(75,0\right)\). 设出发 \(t\mathrm{~h}\) 后两车的位置分别为 \(C\)、\(D\)，其中 \(t\geqslant0\). 由速度和方向，汽车位置为 \(C=\left(75-30t,0\right)\)，自行车位置为 \(D=\left(0,-15t\right)\). 所以两车距离的平方为
\[
CD^2=\left(75-30t\right)^2+\left(15t\right)^2
=1125t^2-4500t+5625
=1125\left(t-2\right)^2+1125
\]
距离本身非负，而平方距离在 \(t\geqslant0\) 上于 \(t=2\) 时取得最小值，因此出发后 \(2\mathrm{~h}\) 两车相距最短.
\end{solution}

\begin{problem}
已知等差数列有 \(10\) 项，偶数项之和为 \(55\)，奇数项之和为 \(45\)，求此数列的第五项.
\end{problem}

\begin{solution}
设等差数列的首项为 \(a_1\)，公差为 \(d\). 其第 \(k\) 项为 \(a_k=a_1+(k-1)d\). 因而奇数项之和为
\[
a_1+a_3+a_5+a_7+a_9=5a_1+\left(0+2+4+6+8\right)d=5a_1+20d=45
\]
偶数项之和为
\[
a_2+a_4+a_6+a_8+a_{10}=5a_1+\left(1+3+5+7+9\right)d=5a_1+25d=55
\]
两式相减，得 \(5d=10\)，所以 \(d=2\). 代入奇数项之和，得 \(5a_1+40=45\)，从而 \(a_1=1\). 因此 \(a_5=a_1+4d=1+4\times2=9\).
\end{solution}

\begin{problem}
一动点到直线 \(x=8\) 的距离是它到点 \(A\left(2,0\right)\) 的距离的 \(2\) 倍.

\begin{enumerate}
\item 求这动点的轨迹方程；
\item 判断这轨迹是何曲线；
\item 求经过这曲线上一点 \(B\left(2,3\right)\) 的切线方程.
\end{enumerate}
\end{problem}

\begin{solution}
设动点为 \(P\left(x,y\right)\).
\begin{enumerate}
\item 点 \(P\) 到直线 \(x=8\) 的距离为 \(\left|x-8\right|\)，到点 \(A\left(2,0\right)\) 的距离为 \(\sqrt{\left(x-2\right)^2+y^2}\). 因此
\[
\left|x-8\right|=2\sqrt{\left(x-2\right)^2+y^2}
\]
两边均非负，平方是等价变形，不引入增根，得
\[
\left(x-8\right)^2=4\left(\left(x-2\right)^2+y^2\right)
\]
展开化简：
\(x^2-16x+64=4x^2-16x+16+4y^2\)，\(3x^2+4y^2=48\).
所以轨迹方程为
\[
\dfrac{x^2}{16}+\dfrac{y^2}{12}=1
\]

\item 将轨迹方程写成标准形式
\[
\dfrac{x^2}{4^2}+\dfrac{y^2}{\left(2\sqrt3\right)^2}=1
\]
所以这条轨迹是以原点为中心、长轴在 \(x\) 轴上的椭圆，半长轴为 \(4\)，半短轴为 \(2\sqrt3\).

\item 先检验给定点：
\[
\dfrac{2^2}{16}+\dfrac{3^2}{12}=\dfrac14+\dfrac34=1
\]
故 \(B\left(2,3\right)\) 在椭圆上. 对等价方程 \(3x^2+4y^2=48\) 隐式求导，得 \(6x+8yy'=0\). 在 \(B\left(2,3\right)\) 处，\(y'=-\dfrac{6\times2}{8\times3}=-\dfrac12\).
因此切线方程为
\(y-3=-\dfrac12\left(x-2\right)\)，\(x+2y=8\).
\end{enumerate}
\end{solution}
